\scanpage{191}
\mspage{II.39}
% archive sheet number: 190
% block: top, lines 1-2 | text
6.9. \uline{Fonctorialité} Un hom de groupes profinis

% block: upper third, line 3 | display
\[
u:\pi\unsure{\text{---}|\longrightarrow}\pi'
\]

% block: upper third, line 4 | text
définit

% block: upper third, lines 5-6 | display
\[
\begin{aligned}
G=\pi^{\mathrm{an}}&\longrightarrow G'=\pi'^{\mathrm{an}},&\operatorname{Lie}(\deleted{\unsure{G}}\pi^{\mathrm{an}})=\mathfrak g&\text{---}\operatorname{Lie}(\pi'^{\mathrm{an}})=\mathfrak g'\\
\widetilde G=\pi^{\mathrm{alg}}&\longrightarrow\widetilde G'=\pi'^{\mathrm{alg}},&\operatorname{Lie}(\deleted{\unsure{G}}\pi^{\mathrm{alg}})=\widetilde{\mathfrak g}&\text{---}\operatorname{Lie}(\pi'^{\mathrm{alg}})=\widetilde{\mathfrak g}'
\end{aligned}
\]

% block: upper third, line 7 | text
avec commutativité dans

% block: upper to middle third, lines 8-9 | diagram
\diagram{\begin{tikzcd}[ampersand replacement=\&]
G=\pi^{\mathrm{an}}\arrow[r]\arrow[d,"u^{\mathrm{an}}"]\&\widetilde G=\pi^{\mathrm{alg}}\arrow[r]\arrow[d,"u^{\mathrm{alg}}"]\&\pi\arrow[d,"u"]\\
G'=\pi'^{\mathrm{an}}\arrow[r,no head]\&\widetilde G'=\pi'^{\mathrm{alg}}\arrow[r,no head]\&\pi'
\end{tikzcd}}

% block: middle third, lines 10-11 | diagram
\diagram{\begin{tikzcd}[ampersand replacement=\&,column sep=large]
\mathfrak g^{\deleted{\unsure{\mathrm{an}}}}=\operatorname{Lie}(\pi^{\mathrm{an}})\arrow[r,no head]\arrow[d]\&\mathfrak g'=\operatorname{Lie}(\pi'^{\mathrm{an}})\arrow[d]\\
\widetilde{\mathfrak g}=\operatorname{Lie}(\pi^{\mathrm{alg}})\arrow[r,no head]\&\widetilde{\mathfrak g}'=\operatorname{Lie}(\pi'^{\mathrm{alg}}).
\end{tikzcd}}

% block: middle third, line 12 | text
\unsure{Transitivités|Transitivité} \unsure{évidentes|évidente}.

% block: middle to lower third, lines 13-17 | text
Si $u:\pi\longrightarrow\pi'$ est un épimorphisme, il en est de même de $u^{\mathrm{an}}$ et $u^{\mathrm{alg}}$ (dans un sens évident). Si l'image de $u$ est d'indice fini, alors $u$ induit un \uline{épimorphisme} $G^0\text{---}G'^0$, et \unsure{des} épimorphismes $\mathfrak g\longrightarrow\mathfrak g'$, $\widetilde{\mathfrak g}\text{---}\widetilde{\mathfrak g}'$.

% block: lower third, lines 18-22 | text
6.10. Soit $\Pi=(\pi_{\unsure{\lambda|\alpha}})$ un système projectif \inserted{filtrant} de groupes profinis. En vertu de 6.9, il définit un système projectif \deleted{filtrant} de systèmes projectifs \deleted{filtrants} de groupes analytiques, groupes algébriques, Alg. de Lie et Alg. de Lie algébriques.

% block: bottom, lines 23-24 | display
\[
\begin{aligned}
G&=\Pi^{\mathrm{an}},&\operatorname{Lie}(\Pi^{\mathrm{an}})&=\mathfrak g\\
\widetilde G&=\Pi^{\mathrm{alg}},&\operatorname{Lie}(\Pi^{\mathrm{alg}})&=\widetilde{\mathfrak g}.
\end{aligned}
\]
